\section{Introduction}\label{sec:introduction}

Let $k$ be a number field and let $G$ be an absolutely almost simple
simply connected algebraic $k$-group. For a place $v$ of $k$, write $k_v$
for the completion. The local group $G(k_v)$ has its usual topology.
Let
\[
 A=A(G,k)=\{v\text{ nonarchimedean}:\rank_{k_v}G=0\},
 \qquad H_A=\prod_{v\in A}G(k_v),
\]
and let $\delta_A:G(k)\to H_A$ be the diagonal homomorphism.
The set $A$ is finite. The product $H_A$ is given its product topology;
if $A$ is empty, it is the trivial group.

The Margulis--Platonov conjecture asserts that the noncentral normal
subgroups of $G(k)$ are accounted for by this finite collection of compact
local groups. Normal subgroups in this assertion are abstract: no topology
or finite-generation condition is imposed on them. We prove the conjecture
in this form.

\begin{theorem}\label{thm:main}
Let $k$ be a number field and $G$ an absolutely almost simple simply
connected algebraic $k$-group. If $N\triangleleft G(k)$ is not contained
in $Z(G(k))$, then there is an open normal subgroup $W\triangleleft H_A$
such that
\[
                         N=\delta_A^{-1}(W).
\]
\end{theorem}

Thus the result resolves the Margulis--Platonov conjecture positively
over number fields, including globally anisotropic groups. In the case
$A=\varnothing$, the conclusion is $N=G(k)$, not merely that $N$ has
finite index or is dense in a local topology.

\subsection{Background and significance}

The problem joins two kinds of rigidity. At isotropic completions,
generation by unipotent subgroups and simplicity modulo the center leave
little room for normal quotients. At anisotropic nonarchimedean
completions, compactness allows finite quotients and their open kernels.
The assertion is that there is no additional abstract normal-subgroup
obstruction invisible to these local factors.

The conjecture grew from Kneser's quaternionic simplicity question and
Platonov's local-to-global simplicity proposal. Margulis formulated the
description by open subgroups of the anisotropic finite local product
\cite[Russian original, Section~2.4.8]{Margulis1979}.
Strong approximation and Margulis's normal subgroup theorem supply the
arithmetic framework and finite-index reduction \cite{PRbook,Margulis}.
The isotropic case follows from the Kneser--Tits theorem over global
fields; Gille's account includes the final outer-$E_6$ cases, using work
of Garibaldi and Chernousov--Timoshenko
\cite[Theorem~8.1]{Gille2009}. For anisotropic groups, Chernousov proved
projective simplicity of the rational points when the absolute rank is
at least two and the group splits over a quadratic extension
\cite[Theorem~1]{Chernousov1990}. Tomanov and Sury established substantial
classical and low-index unitary cases \cite{Tomanov1990,Sury1991}.

For inner type $A$, the arithmetic work of Platonov--Rapinchuk and
Raghunathan \cite{PR1984,Raghunathan1988} preceded the reduction of
Rapinchuk--Potapchik to finite simple quotients of the multiplicative
group of a division algebra \cite{RapinchukPotapchik1996}.
Segev obtained a commuting-graph obstruction, and Segev--Seitz verified
the required finite-simple-group properties, completing this case
\cite{Segev1999,SegevSeitz2002}. Rapinchuk--Segev--Seitz later proved
that every finite quotient of the multiplicative group of a
finite-dimensional division algebra is solvable
\cite{RapinchukSegevSeitz2002}; Rapinchuk gave a shorter proof of the
inner-$A$ case using two-generation of finite simple groups
\cite{Rapinchuk2006}. The established cases are surveyed in
\cite[Sections~3.4--3.5]{PRsurvey}. They include isotropic groups and
all inner forms of type $A$, as well as the other types outside the
remaining anisotropic outer $A$, trialitarian $D_4$, and $E_6$ cases.
We use that reduction in Theorem~\ref{thm:known-mp} and supply the
remaining arguments here. In particular, the previously established
inner-type-$A$ theorem is available over every number field that occurs
in a restriction-of-scalars subgroup.

The conjecture also enters the congruence subgroup problem: it is a
normal-subgroup input to several descriptions of congruence kernels.
That relationship does not justify using a theorem which assumes the
conjecture to prove it. Our central-extension argument instead uses the
unconditional finite-nonarchimedean injectivity in the metaplectic-kernel
computation of Prasad and Rapinchuk \cite[Proposition~2.6]{PRmetaplectic}.
The distinction between metaplectic vanishing and congruence-kernel
conclusions under the normal-subgroup hypothesis is also explicit in
\cite[Section~2]{Rapinchuk2026}.

\subsection{A positive-density congruence consequence}

Theorem~\ref{thm:main} supplies the Margulis--Platonov hypothesis in
Rapinchuk's positive-density theorem, giving the following consequence.

\begin{corollary}[Positive-density congruence subgroup property]
\label{cor:positive-density-csp}
Let $k$ be a number field and $G$ an absolutely almost simple simply
connected algebraic $k$-group. Let $M/k$ be the minimal Galois extension
over which $G$ becomes an inner form of the split group, and let
$\operatorname{Spl}(M/k)$ be the set of nonarchimedean places of $k$ that
split completely in $M$. Suppose that $S$ is a set of places of $k$ which
contains all archimedean places, contains no nonarchimedean place $v$
with $\rank_{k_v}G=0$, and satisfies
\[
 d_k\bigl(S\cap\operatorname{Spl}(M/k)\bigr)>0,
\]
where $d_k$ denotes Dirichlet density and the indicated density is assumed
to exist. Then the $S$-congruence kernel is trivial:
\[
 C^S(G)=\{1\}.
\]
Equivalently, for any fixed faithful $k$-defined representation of $G$,
the associated group $G(\mathcal O(S))$ over the ring of $S$-integers
has the classical congruence subgroup property: every finite-index normal
subgroup contains a principal congruence subgroup of some nonzero ideal
level.
\end{corollary}

\begin{proof}
Theorem~\ref{thm:main} is the Margulis--Platonov statement assumed in
\cite[Theorem~1.1]{Rapinchuk2026}, for the same group and number field.
The other hypotheses are exactly those imposed above, so that theorem
gives $C^S(G)=\{1\}$. The equivalence with the classical congruence
subgroup property is recalled in \cite[Section~2]{Rapinchuk2026}.
\end{proof}

For an inner form of the split group, $M=k$, and the density condition is
positive Dirichlet density of the nonarchimedean places in $S$. The
corollary is restricted to the stated positive-density range: it does not
address Serre's finite-$S$ conjecture, arbitrary infinite $S$, or
centrality or finiteness of congruence kernels outside that range.

\subsection{The finite obstruction}

We first prove that a noncentral normal subgroup $N$ has finite index.
Choose a positive integer $e$ that annihilates the finite group $G(k)/N$,
and let
\[
 R=G(k)^e=\langle g^e:g\in G(k)\rangle,\qquad
 P_0=\overline{\delta_A(R)},\qquad
 V_0=\delta_A^{-1}(P_0).
\]
Here $G(k)^e$ is the subgroup generated by powers, not just the set of
powers. The group $P_0$ is open and normal in $H_A$. It remains to show
that the finite group $V_0/R$ is trivial. This is the distinction between
the desired abstract equality and the closure statement furnished by
approximation.

Section~\ref{sec:foundations} develops two tools for this finite
obstruction. First, it constructs representatives of a fixed coset
modulo $R$ satisfying specified local openings and avoiding a prescribed
codimension-two reduction locus at almost all other places. The parameters
are products of conjugates of torus powers. Scalar invariance permits an
affine sieve while retaining control at the omitted approximation place;
possible single poles are treated explicitly. Second, a signed permutation
lattice calculation gives the Hasse principle and weak approximation for
the norm torus needed to prescribe commuting elements with exact
determinants.

\subsection{Odd-degree division algebras}

The first unresolved case is a special unitary group $S=\SU(D,*)$, where
$D$ is an odd-degree central division algebra over a quadratic extension
$L/k$ and $*$ is a unitary involution. Let $U=\U(D,*)$ be its full
unitary group. Section~\ref{sec:colors} proves that a nontrivial finite
defect for $S(k)$ yields one of two objects on $U(k)$: a finite abelian
character nonzero on $S(k)$, or a homomorphism to a finite nonabelian
simple group which is surjective on the defect preimage $V_0$.
The commuting approximation and central-extension arguments are needed
to obtain this dichotomy on the full unitary group. The class-preserving
extension step adapts the inner-type argument of Rapinchuk--Potapchik
\cite[Section~2]{RapinchukPotapchik1996}; here the commuting approximation
must also prescribe determinants. A quotient of a unitary subgroup need
not extend to the multiplicative group of its algebra.

Section~\ref{sec:characters} rules out the first object. An algebraic
difference function on $D$ has enough exact elementary invariances to be
constant on the required orbits. The exact additive span of $U(k)$ in
$D$, together with the known inner-type-$A$ theorem and a real-signature
interpolation argument, forces any finite abelian character to vanish
on $S(k)$.

Sections~\ref{sec:palettes} and~\ref{sec:finite-groups} rule out the
second object. A finite simple quotient assigns a color to each rational
unitary element. The distribution of these colors is not assumed
continuous. We construct a translation-invariant probability law on
their finite patterns whose individual color distributions are uniform.
Fractional unitary transformations turn additive recurrence into a
centralizer double-coset condition. A finite-simple-group lemma, stated
in Section~\ref{sec:palettes} and proved in
Section~\ref{sec:finite-groups}, excludes that condition. Elementary
transformations then force an impossible invariance under all left
color translations.

The construction of the uniform-marginal law is an important part of the
proof: additive averaging by itself need not preserve all colors.
Conformal weights on rational points of a flag variety use Langlands's
Eisenstein-series convergence theorem \cite{Langlands}. Finite-volume
harmonicity and Howe--Moore matrix-coefficient decay
\cite{HoweMoore,Ciobotaru} then provide the required marginals before
additive averaging is performed. The later positive-density argument
uses the multidimensional theorem of Furstenberg--Katznelson
\cite{FurstenbergKatznelson}. The uniform-marginal construction, the
coset-preserving sieve, and the explicit norm-torus calculation are
developed before the final case reductions so that their hypotheses and
conclusions remain visible.

\subsection{Completion of the proof}

Section~\ref{sec:unitary-induction} proves the other outer-type-$A$
cases by induction on reduced degree, simultaneously over all number
fields. The ternary split-algebra case uses quadratic norm equations on a
Ch\^atelet surface. Even degree is reduced by a quaternion factorization
and a smaller unitary centralizer, with a separate degree-four determinant
correction. The approximation construction controls the finite places of
the subgroups as they vary.

\enlargethispage{\baselineskip}
Section~\ref{sec:d4} treats $D_4$ by finding, in each relevant rational
coset, a regular element whose torus admits a rational inverter. It is a
product of two involutions, each already controlled inside a type-$A_2$
subgroup. Section~\ref{sec:e6} reduces a generic $E_6$ element to a
simply connected $D_4$ subgroup and the cover of a root-involution
centralizer of type $A_1A_5$. The last step is a torsor under a connected
simply connected group, so its local factorizations can be globalized.
The proof ends by returning to the original abstract subgroup $N$.
